\documentclass[../../main.tex]{subfiles}% 注意这里的文件路径不能用 ./main.tex ,否则用latexmk编译子文件会报错
\graphicspath{{\subfix{./image/}}} % 指定图片目录,后续可以直接使用图片文件名
% 注意这里的文件路径不能用 ../../image/ ,否则用latexmk编译子文件会报错

% 例如:
% \begin{figure}[H]
% \centering
% \includegraphics[scale=0.3]{图.png}
% \caption{}
% \label{figure:图}
% \end{figure}
% 注意:上述\label{}一定要放在\caption{}之后,否则引用图片序号会只会显示??.

\begin{document}

\section{Sobolev空间与嵌入定理}

\begin{definition}
设$\Omega\subseteq\mathbb{R}^{n}$为开集，$m$为非负整数，$1\leqslant p<\infty$，定义Sobolev空间
\begin{align*}
W^{m,p}(\Omega)=\{u\in L^{p}(\Omega)\mid \widetilde{\partial}^{\alpha}u\in L^{p}(\Omega),|\alpha|\leqslant m\}.
\end{align*}
其上赋予范数
\begin{align*}
\|u\|_{m,p}=\left(\sum\limits_{|\alpha|\leqslant m}\|\widetilde{\partial}^{\alpha}u\|_{L^{p}(\Omega)}^{p}\right)^{\frac{1}{p}}=\left(\sum\limits_{|\alpha|\leqslant m}\int_{\Omega}|\widetilde{\partial}^{\alpha}u(x)|^{p}\mathrm{d}x\right)^{\frac{1}{p}},
\end{align*}
该赋范空间称为\textbf{Sobolev空间}，也可记作$W_{p}^{m}(\Omega)$。
\end{definition}

\begin{theorem}
空间$W^{m,p}(\Omega)$是完备的。
\end{theorem}

\begin{proof}
设$\{u_{k}\}_{k=1}^{\infty}$为$W^{m,p}(\Omega)$中的基本列，则对任意$|\alpha|\leqslant m$，$\{\widetilde{\partial}^{\alpha}u_{k}\}$是$L^{p}(\Omega)$中的基本列。
由$L^{p}(\Omega)$的完备性，存在$g_{\alpha}\in L^{p}(\Omega)$，使得
\begin{align*}
\|\widetilde{\partial}^{\alpha}u_{k}-g_{\alpha}\|_{L^{p}(\Omega)}\to 0,\quad k\to\infty,\ |\alpha|\leqslant m.
\end{align*}
对任意$\varphi\in\mathscr{D}(\Omega)$，当$k\to\infty$时，
\begin{align*}
\langle\widetilde{\partial}^{\alpha}u_{k},\varphi\rangle=(-1)^{|\alpha|}\langle u_{k},\partial^{\alpha}\varphi\rangle,
\end{align*}
两边取极限得$\langle g_{\alpha},\varphi\rangle=(-1)^{|\alpha|}\langle g_{0},\partial^{\alpha}\varphi\rangle$，即$g_{\alpha}=\widetilde{\partial}^{\alpha}g_{0}$。
因此$g_{0}\in W^{m,p}(\Omega)$，且$\|u_{k}-g_{0}\|_{m,p}\to 0,\ k\to\infty$.

\end{proof}

\begin{theorem}
$\mathscr{D}(\mathbb{R}^{n})$在$W^{m,p}(\mathbb{R}^{n})$中稠密。
\end{theorem}

\begin{proof}
任取$u\in W^{m,p}(\mathbb{R}^{n})$，则$\widetilde{\partial}^{\alpha}u\in L^{p}(\mathbb{R}^{n}),\ |\alpha|\leqslant m$。对任意$\delta>0$，定义磨光函数
\begin{align}
u_{\delta}(x)=\int_{\mathbb{R}^{n}}u(y)\cdot j_{\delta}(x-y)\mathrm{d}y,\label{eq:mollify}
\end{align}
其中$j_{\delta}$为标准磨光核。可验证$u_{\delta}\in C^{\infty}(\mathbb{R}^{n})$，且
\begin{align*}
\partial_{x}^{\alpha}u_{\delta}(x)=\int_{\mathbb{R}^{n}}u(y)\cdot\partial_{x}^{\alpha}j_{\delta}(x-y)\mathrm{d}y=\int_{\mathbb{R}^{n}}\widetilde{\partial}^{\alpha}u(y)\cdot j_{\delta}(x-y)\mathrm{d}y=(\widetilde{\partial}^{\alpha}u)_{\delta}.
\end{align*}
由Young不等式，对任意$v\in L^{p}(\mathbb{R}^{n})$，有
\begin{align}
\|v_{\delta}\|_{L^{p}}\leqslant\|v\|_{L^{p}}.\label{eq:young}
\end{align}
对$\widetilde{\partial}^{\alpha}u\in L^{p}(\mathbb{R}^{n})$，对任意$\varepsilon>0$，存在$v_{\alpha}\in C_{0}^{\infty}(\mathbb{R}^{n})$，使得
\begin{align}
\|v_{\alpha}-\widetilde{\partial}^{\alpha}u\|_{L^{p}}<\frac{\varepsilon}{3}.\label{eq:approx1}
\end{align}
结合\eqref{eq:young}得
\begin{align}
\|(v_{\alpha})_{\delta}-(\widetilde{\partial}^{\alpha}u)_{\delta}\|_{L^{p}}<\frac{\varepsilon}{3}.\label{eq:approx2}
\end{align}
由磨光函数逼近性质，存在$\delta_{0}=\delta_{0}(\varepsilon,\alpha)>0$，当$0<\delta<\delta_{0}$时，
\begin{align}
\|(v_{\alpha})_{\delta}-v_{\alpha}\|_{L^{p}}<\frac{\varepsilon}{3}.\label{eq:approx3}
\end{align}
联立\eqref{eq:approx1},\eqref{eq:approx2},\eqref{eq:approx3}，得$\|(\widetilde{\partial}^{\alpha}u)_{\delta}-\widetilde{\partial}^{\alpha}u\|_{L^{p}}<\varepsilon$，即$\|\partial^{\alpha}u_{\delta}-\widetilde{\partial}^{\alpha}u\|_{L^{p}}<\varepsilon,\ |\alpha|\leqslant m$.

\end{proof}

\begin{remark}
记$W_{0}^{m,p}(\mathbb{R}^{n})$为$C_{0}^{\infty}(\mathbb{R}^{n})$按范数$\|\cdot\|_{m,p}$完备化生成的空间，则$W_{0}^{m,p}(\mathbb{R}^{n})=W^{m,p}(\mathbb{R}^{n})$；对一般开集$\Omega$，该等式未必成立。
在第一章中定义$H^{m,p}(\Omega)$为集合$S=\{u\in C^{\infty}(\Omega)\mid\|u\|_{m,p}<\infty\}$在范数$\|\cdot\|_{m,p}$下的完备化空间，由$W^{m,p}(\Omega)$的完备性，有$H^{m,p}(\Omega)\subseteq W^{m,p}(\Omega)$，二者实际等价，需用到$C^{\infty}$单位分解定理。
\end{remark}

\begin{theorem}
设$A\subseteq\mathbb{R}^{n}$，$\mathscr{O}$为$A$的开覆盖，则存在一族$C^{\infty}$函数$\mathscr{F}$，满足：
\begin{enumerate}[(1)]
\item $0\leqslant\varphi(x)\leqslant 1,\ \forall x\in A$；
\item 对任意$x\in A$，存在含$x$的开集$V$，使得仅有有限个$\varphi\in\mathscr{F}$在$V$上非零；
\item $\sum\limits_{\varphi\in\mathscr{F}}\varphi(x)\equiv 1,\ \forall x\in A$；
\item 对任意$\varphi\in\mathscr{F}$，存在$U\in\mathscr{O}$，使得$\mathrm{supp}(\varphi)\subseteq U$。
\end{enumerate}
\end{theorem}

\begin{theorem}[Meyers‑Serrin]
若$1\leqslant p<\infty$，则$H^{m,p}(\Omega)=W^{m,p}(\Omega)$。
\end{theorem}

\begin{proof}
只需证$S=\{u\in C^{\infty}(\Omega)\mid\|u\|_{m,p}<\infty\}$在$W^{m,p}(\Omega)$中稠密。
记
\begin{align*}
\Omega_{k}=\left\{x\in\Omega\mid\|x\|<k,\ \mathrm{dist}(x,\partial\Omega)>\frac{1}{k}\right\},\quad k=1,2,\dots,
\end{align*}
$\Omega_{0}=\Omega_{-1}=\emptyset$，取开集族$\mathscr{O}=\{U_{k}\mid U_{k}=\Omega_{k+1}\cap C\overline{\Omega}_{k-1},\ k=1,2,\dots\}$，其为$\Omega$的开覆盖。
取$\mathscr{O}$对应的$C^{\infty}$单位分解$\mathscr{F}$，记$\psi_{k}$为$\mathscr{F}$中支集含于$U_{k}$的函数之和，则$\psi_{k}\in C_{0}^{\infty}(U_{k})$，且$\sum\limits_{k=1}^{\infty}\psi_{k}(x)\equiv 1,\ \forall x\in\Omega$。
任取$u\in W^{m,p}(\Omega)$，$\forall\varepsilon>0$，取$0<\delta<\frac{1}{(k+1)(k+2)}$，则$\mathrm{supp}(\psi_{k}u)_{\delta}\subseteq\Omega_{k+2}\cap C\overline{\Omega}_{k-2}$，存在$\delta_{k}\in\left(0,\frac{1}{(k+1)(k+2)}\right)$，使得
\begin{align*}
\|\psi_{k}u-(\psi_{k}u)_{\delta_{k}}\|_{m,p}<\frac{\varepsilon}{2^{k}}.
\end{align*}
令$\varphi=\sum\limits_{k=1}^{\infty}(\psi_{k}u)_{\delta_{k}}$，因任意点$x\in\Omega$的邻域内仅有有限项非零，故$\varphi\in C^{\infty}(\Omega)$。
在$\Omega_{k}$上，$u=\sum\limits_{j=1}^{k+2}\psi_{j}u$，$\varphi=\sum\limits_{j=1}^{k+2}(\psi_{j}u)_{\delta_{j}}$，因此
\begin{align*}
\|u-\varphi\|_{W^{m,p}(\Omega_{k})}\leqslant\sum\limits_{j=1}^{k+2}\|(\psi_{j}u)_{\delta_{j}}-\psi_{j}u\|_{m,p}<\varepsilon.
\end{align*}
令$k\to\infty$，得$\|u-\varphi\|_{m,p}\leqslant\varepsilon$，且$\varphi\in S$，故稠密性成立.

\end{proof}

\begin{definition}
$\mathbb{R}^{n}$中的开区域$\Omega$称为\textbf{可扩张的}，若对任意$m\in\mathbb{N}$，$p\in[1,\infty)$，存在连续线性算子$T:W^{m,p}(\Omega)\to W^{m,p}(\mathbb{R}^{n})$，满足$Tu|_{\Omega}=u,\ \forall u\in W^{m,p}(\Omega)$。
\end{definition}

\begin{example}
$\mathbb{R}_{+}^{n}=\{(x_{1},x_{2},\dots,x_{n})\in\mathbb{R}^{n}\mid x_{n}>0\}$是可扩张的。
\end{example}

\begin{proof}
对任意$m\in\mathbb{N}$，$u\in C^{\infty}(\mathbb{R}_{+}^{n})$，定义扩张算子$E_{m}$：
\begin{align*}
E_{m}u(x)=
\begin{cases}
u(x),&x_{n}>0,\\
\sum\limits_{j=1}^{m+1}\lambda_{j}u(x_{1},\dots,x_{n-1},-jx_{n}),&x_{n}\leqslant 0,
\end{cases}
\end{align*}
其中系数$\lambda_{1},\lambda_{2},\dots,\lambda_{m+1}$是线性方程组
\begin{align*}
\sum\limits_{j=1}^{m+1}(-j)^{i}\lambda_{j}=1,\quad i=0,1,\dots,m
\end{align*}
的唯一解。
若$u\in C^{m}(\overline{\mathbb{R}_{+}^{n}})$，则$E_{m}u\in C^{m}(\mathbb{R}^{n})$，且
\begin{align*}
\|E_{m}u\|_{W^{m,p}(\mathbb{R}^{n})}\leqslant M_{m,p}\|u\|_{W^{m,p}(\mathbb{R}_{+}^{n})},
\end{align*}
$M_{m,p}$为常数。$E_{m}$可连续延拓至$W^{m,p}(\mathbb{R}_{+}^{n})$，故$\mathbb{R}_{+}^{n}$可扩张.

\end{proof}

\begin{example}
若$\Omega$为有界开区域，具有一致$C^{m}$光滑边界，则$\Omega$是可扩张的。
\end{example}

\begin{theorem}[Sobolev嵌入定理]
若$\Omega\subseteq\mathbb{R}^{n}$为可扩张区域，$m>\frac{n}{2}$，则$W^{m,2}(\Omega)$可连续嵌入$C(\overline{\Omega})$。
\end{theorem}

\begin{proof}
\begin{enumerate}[(1)]
\item 取$W^{m,2}(\mathbb{R}^{n})$的等价范数
\begin{align*}
\|u\|_{m}'=\left(\int_{\mathbb{R}^{n}}(1+|\xi|^{2})^{m}|\widehat{u}(\xi)|^{2}\mathrm{d}\xi\right)^{\frac{1}{2}}.
\end{align*}
$L^{1}(\mathbb{R}^{n})$函数的Fourier逆变换为连续函数，且
\begin{align*}
\|u\|_{C(\mathbb{R}^{n})}\leqslant\int_{\mathbb{R}^{n}}|\widehat{u}(\xi)|\mathrm{d}\xi.
\end{align*}
当$m>\frac{n}{2}$时，由Cauchy‑Schwarz不等式：
\begin{align*}
\int_{\mathbb{R}^{n}}|\widehat{u}(\xi)|\mathrm{d}\xi&\leqslant\left(\int_{\mathbb{R}^{n}}(1+|\xi|^{2})^{m}|\widehat{u}(\xi)|^{2}\mathrm{d}\xi\right)^{\frac{1}{2}}\left(\int_{\mathbb{R}^{n}}\frac{\mathrm{d}\xi}{(1+|\xi|^{2})^{m}}\right)^{\frac{1}{2}}\\
&\leqslant c_{n,m}\|u\|_{m}',
\end{align*}
故$i:u\mapsto u$是$W^{m,2}(\mathbb{R}^{n})\to C(\mathbb{R}^{n})$的连续嵌入。
\item 对任意$u\in W^{m,2}(\Omega)$，取扩张算子$T$，则
\begin{align*}
\|u\|_{C(\overline{\Omega})}&\leqslant\|Tu\|_{C(\mathbb{R}^{n})}\leqslant c_{n,m}\|Tu\|_{m}'\\
&\leqslant c\|u\|_{W^{m,2}(\Omega)},
\end{align*}
故$i:u\mapsto u$是$W^{m,2}(\Omega)\to C(\overline{\Omega})$的连续嵌入。
\end{enumerate}

\end{proof}

\begin{remark}
一般Sobolev嵌入结论：
\begin{align*}
W^{m,p}(\Omega)\hookrightarrow L^{q}(\Omega),\quad \frac{1}{q}=\frac{1}{p}-\frac{m}{n},\ m\leqslant\frac{n}{p},\\
W^{m,p}(\Omega)\hookrightarrow C(\overline{\Omega}),\quad m>\frac{n}{p},
\end{align*}
其中$\hookrightarrow$表示连续嵌入。
\end{remark}

\begin{theorem}[Rellich定理]
若$\Omega\subseteq\mathbb{R}^{n}$为有界可扩张区域，则$W^{1,2}(\Omega)$的单位球在$L^{2}(\Omega)$中列紧。
\end{theorem}

\begin{proof}
设$\{u_{m}\}_{m=1}^{\infty}\subseteq W^{1,2}(\Omega)$，$\|u_{m}\|_{W^{1,2}}\leqslant 1$，取扩张算子$T$，记$U_{m}=Tu_{m}$，不妨设$U_{m}$在紧集$K\subseteq\mathbb{R}^{n}$外为0。
只需证$\{U_{m}\}$存在子列在$L^{2}(\mathbb{R}^{n})$中收敛。
计算得
\begin{align*}
\|u_{p'}-u_{m'}\|_{L^{2}(\Omega)}^{2}\leqslant C\|U_{p'}-U_{m'}\|_{L^{2}(\mathbb{R}^{n})}^{2}=C\int_{\mathbb{R}^{n}}|\widehat{U}_{p'}(\xi)-\widehat{U}_{m'}(\xi)|^{2}\mathrm{d}\xi,
\end{align*}
其中$\widehat{U}_{m}(\xi)=\int_{K}U_{m}(x)e^{-2\pi i x\cdot\xi}\mathrm{d}x=\langle e^{-2\pi i x\cdot\xi}\alpha_{K},U_{m}\rangle$，$\alpha_{K}(x)=\begin{cases}1,&x\in K,\\0,&x\notin K.\end{cases}$
由$L^{2}(\mathbb{R}^{n})$的自反性与Eberlein‑Smulian定理，存在子列$\{U_{m'}\}$弱收敛，故$\{\widehat{U}_{m'}(\xi)\}$逐点收敛，且$|\widehat{U}_{m}(\xi)|\leqslant[\mathrm{mes}(K)]^{\frac{1}{2}}\|U_{m}\|_{L^{2}}\leqslant\mathrm{const}$。
对任意$r>0$，拆分积分
\begin{align*}
\int_{\mathbb{R}^{n}}|\widehat{U}_{m'}(\xi)-\widehat{U}_{p'}(\xi)|^{2}\mathrm{d}\xi=\int_{|\xi|\leqslant r}+\int_{|\xi|\geqslant r}=I+\mathrm{II}.
\end{align*}
估计第二项：
\begin{align*}
\mathrm{II}&\leqslant r^{-2}\int_{|\xi|\geqslant r}(1+|\xi|^{2})|\widehat{U}_{m'}(\xi)-\widehat{U}_{p'}(\xi)|^{2}\mathrm{d}\xi\\
&\leqslant 2r^{-2}\int_{|\xi|\geqslant r}(1+|\xi|^{2})\left(|\widehat{U}_{m'}(\xi)|^{2}+|\widehat{U}_{p'}(\xi)|^{2}\right)\mathrm{d}\xi\\
&\leqslant 2r^{-2}\left(\|U_{m'}\|_{1}'^{2}+\|U_{p'}\|_{1}'^{2}\right)\leqslant Mr^{-2},
\end{align*}
$M$为仅依赖扩张算子的常数。
对任意$\varepsilon>0$，取$r$充分大使得$Mr^{-2}<\frac{\varepsilon}{2}$，固定$r$，由Lebesgue控制收敛定理，当$m',p'$充分大时，$I<\frac{\varepsilon}{2}$，故$\|u_{m'}-u_{p'}\|_{L^{2}(\Omega)}\to 0,\ m',p'\to\infty$.

\end{proof}

\begin{corollary}
若$\Omega\subseteq\mathbb{R}^{n}$为任意有界开集，则$W_{0}^{1,2}(\Omega)$的单位球在$L^{2}(\Omega)$中列紧。
\end{corollary}
\begin{remark}
\begin{enumerate}[(1)]
\item $W_{0}^{1,2}(\Omega)$也记作$\mathring{W}_{2}^{1}(\Omega)$，由Meyers‑Serrin定理，其即为$H_{0}^{1}(\Omega)$。

\item 一般结论由Konttrashev给出，可参考Adams的\textit{Sobolev Spaces}中定理6.2。

\item 由$\mathscr{D}(\Omega)\hookrightarrow H_{0}^{m}(\Omega)\hookrightarrow L^{2}(\Omega)$，其共轭空间满足$L^{2}(\Omega)\hookrightarrow H_{0}^{m}(\Omega)^{*}\hookrightarrow\mathscr{D}'(\Omega)$，故$H_{0}^{m}(\Omega)^{*}$包含比$L^{2}(\Omega)$更多的广义函数，记$H_{0}^{m}(\Omega)^{*}=H^{-m}(\Omega)$。
\end{enumerate}
\end{remark}
\begin{proof}
证明过程与Rellich定理类似.

\end{proof}



\begin{theorem}
$f\in H^{-m}(\Omega)$当且仅当存在$g_{\alpha}\in L^{2}(\Omega),\ |\alpha|\leqslant m$，使得
\begin{align*}
\langle f,\varphi\rangle=\sum\limits_{|\alpha|\leqslant m}\int_{\Omega}g_{\alpha}(x)\cdot\partial^{\alpha}\varphi(x)\mathrm{d}x,\quad \forall\varphi\in H_{0}^{m}(\Omega).
\end{align*}
\end{theorem}

\begin{proof}
充分性显然，只需证必要性。
任取$f\in H^{-m}(\Omega)$，因$H_{0}^{m}(\Omega)$为Hilbert空间，由Riesz表示定理，存在$h\in H_{0}^{m}(\Omega)$，使得
\begin{align}
\langle f,\varphi\rangle=(\varphi,h)_{m}=\sum\limits_{|\alpha|\leqslant m}\int_{\Omega}\overline{\widetilde{\partial}^{\alpha}h(x)}\cdot\widetilde{\partial}^{\alpha}\varphi(x)\mathrm{d}x=\sum\limits_{|\alpha|\leqslant m}\int_{\Omega}g_{\alpha}(x)\widetilde{\partial}^{\alpha}\varphi(x)\mathrm{d}x,\quad \forall\varphi\in H_{0}^{m}(\Omega),\label{eq:riesz111}
\end{align}
其中$g_{\alpha}(x)=\widetilde{\partial}^{\alpha}h(x)\in L^{2}(\Omega)$.

\end{proof}

\begin{corollary}
每个$f\in H^{-m}(\Omega)$是$L^{2}(\Omega)$函数广义微商的有限和，即
\begin{align}
f=(-1)^{|\alpha|}\sum\limits_{|\alpha|\leqslant m}\widetilde{\partial}^{\alpha}g_{\alpha},\quad g_{\alpha}\in L^{2}(\Omega).\label{eq:negorder}
\end{align}
\end{corollary}
\begin{remark}
\begin{enumerate}[(1)]
\item 给定$g_{\alpha}\in L^{2}(\Omega)$，由\eqref{eq:negorder}定义的$f\in\mathscr{D}'(\Omega)$在$H^{m}(\Omega)$上的连续延拓不唯一，但在$H_{0}^{m}(\Omega)$上唯一，因$C_{0}^{\infty}(\Omega)$在$H_{0}^{m}(\Omega)$中稠密，且$|\langle\widetilde{\partial}^{\alpha}g_{\alpha},\varphi\rangle|\leqslant\|g_{\alpha}\|_{L^{2}}\|\varphi\|_{H_{0}^{m}(\Omega)},\ \forall\varphi\in C_{0}^{\infty}(\Omega)$。
\item 由Sobolev嵌入定理，当$m>\frac{n}{2}$时，$H_{0}^{m}(\Omega)\hookrightarrow C(\overline{\Omega})$，故$C(\overline{\Omega})$上的连续线性泛函属于$H^{-m}(\Omega)$，特别地$\delta$函数属于$H^{-m}(\Omega)$。例如$n=1$，$\Omega=(a,b)$，$\varphi\mapsto\langle\delta_{x_{0}},\varphi\rangle=\varphi(x_{0})$在$C(\overline{\Omega})$上连续，故$\delta_{x_{0}}\in H^{-1}(\Omega)$。
\end{enumerate}
\end{remark}
\begin{proof}
由广义微商定义，式\eqref{eq:riesz111}蕴含式\eqref{eq:negorder}.

\end{proof}

\begin{example}
就 $\Omega=\mathbb{R}_+^n=\{(x_1,x_2,\cdots,x_n)\in\mathbb{R}^n\mid x_n>0\}$ 的情形验证定理4.5.5.

提示 $\forall u(x)\in W^{m,p}(\Omega),\forall\varepsilon>0$，请考虑 $u_\varepsilon(x)=u(x',x_n+\varepsilon)$，其中 $x'\in\mathbb{R}^{n-1},x_n>-\varepsilon$.
\end{example}

\begin{example}
若 $\alpha\in\mathscr{D},u\in W^{m,p}(\mathbb{R}^n)$，则 $\alpha\cdot u\in W^{m,p}(\mathbb{R}^n)$，并且有常数 $C$（依赖于 $\alpha$），使得
\begin{align*}
\|\alpha\cdot u\|_{W^{m,p}}\leqslant C\|u\|_{W^{m,p}}.
\end{align*}
\end{example}

\begin{example}
若 $m\geqslant l$，求证：$W^{m,p}(\Omega)\hookrightarrow W^{l,p}(\Omega)$.
\end{example}

\begin{example}
设 $\Omega=(a,b),\forall f\in L^2(\Omega)$，求证：$\exists x\in H_0^1(\Omega)$，使得
\begin{align*}
\frac{\widetilde{\mathrm{d}}^2x}{\mathrm{d}t^2}=f,
\end{align*}
并且 $T:f\mapsto x$ 是 $L^2(\Omega)$ 到 $H^2(\Omega)$ 的连续线性算子.
\end{example}

\begin{example}
设 $f(x)\in H_0^1(-1,1)$，求证：
\begin{enumerate}[(1)]
\item $f(-1)=f(1)=0$；
\item $f(x)$ 绝对连续；
\item $f'(x)\in L^2(-1,1)$（这里$'$指的是求a.e.微商）.
\end{enumerate}
\end{example}

\begin{example}
设 $f\in H^s(\mathbb{R}^n)$（定义见习题4.4.2），求证：当 $s>\frac{n}{2}$ 时，
\begin{enumerate}[(1)]
\item $\widehat{f}(\xi)\in L^1(\mathbb{R}^n)$；
\item $f(x)$ 与一个 $\mathbb{R}^n$ 上的连续有界函数几乎处处相等.
\end{enumerate}
\end{example}

\begin{example}
设 $m\in\mathbb{N}$，又设
\begin{align*}
H^{-m}\triangleq\left\{f\in\mathscr{S}'\,\bigg|\,(1+|\xi|^2)^{-\frac{m}{2}}\widehat{f}(\xi)\in L^2(\mathbb{R}^n)\right\},
\end{align*}
在 $H^{-m}$ 中定义范数
\begin{align*}
\|f\|_{-m}=\left\|(1+|\xi|^2)^{-\frac{m}{2}}\widehat{f}(\xi)\right\|_{L^2(\mathbb{R}^n)}\quad(\forall f\in H^{-m}).
\end{align*}
求证：若 $f\in H^{-m}$，则它可以表为有限个 $L^2(\mathbb{R}^n)$ 函数的导数之和.

提示 $(1+|\xi|^2)^{-\frac{m}{2}}\widehat{f}(\xi)\in L^2(\mathbb{R}^n)$
\begin{align*}
\iff \frac{\widehat{f}(\xi)}{1+|\xi_1|^m+\cdots+|\xi_n|^m}\in L^2(\mathbb{R}^n).
\end{align*}
\end{example}

\begin{example}
在空间 $L^2(-\infty,\infty)$ 上，考察微分算子
\begin{align*}
A=\frac{\widetilde{\mathrm{d}}}{\mathrm{d}x},\quad D(A)=H^1(-\infty,\infty),
\end{align*}
求证：
\begin{enumerate}[(1)]
\item $\rho(A)=\{\lambda\in\mathbb{C}\mid\mathrm{Re}\,\lambda\neq0\}$；
\item $\sigma_p(A)=\emptyset$；
\item $\sigma(A)=\sigma_c(A)=\{\lambda\in\mathbb{C}\mid\mathrm{Re}\,\lambda=0\}$.
\end{enumerate}
\end{example}





















\end{document}